\documentclass[11pt,a4paper]{article}

%% ── Packages ──────────────────────────────────────────────────────────────
\usepackage[utf8]{inputenc}
\usepackage[T1]{fontenc}
\usepackage{amsmath,amssymb,amsthm}
\usepackage{algorithm}
\usepackage{algpseudocode}
\usepackage{geometry}
\usepackage{hyperref}
\usepackage{booktabs}
\usepackage{enumitem}
\usepackage{tikz}
\usepackage{graphicx}
\usepackage{xcolor}
\usepackage{fancyhdr}
\usepackage{titlesec}
\usepackage{tocloft}
\usepackage{multicol}

\geometry{margin=1in}
\hypersetup{
  colorlinks=true,
  linkcolor=blue!70!black,
  citecolor=green!50!black,
  urlcolor=blue!60!black
}

%% ── Theorem environments ──────────────────────────────────────────────────
\theoremstyle{definition}
\newtheorem{definition}{Definition}[section]
\newtheorem{theorem}{Theorem}[section]
\newtheorem{lemma}[theorem]{Lemma}
\newtheorem{proposition}[theorem]{Proposition}
\newtheorem{corollary}[theorem]{Corollary}
\newtheorem{remark}{Remark}[section]
\newtheorem{example}{Example}[section]

%% ── Custom commands ───────────────────────────────────────────────────────
\newcommand{\R}{\mathbb{R}}
\newcommand{\N}{\mathbb{N}}
\newcommand{\calO}{\mathcal{O}}
\newcommand{\norm}[1]{\left\lVert #1 \right\rVert}
\newcommand{\inner}[2]{\left\langle #1, #2 \right\rangle}
\newcommand{\dist}{\mathrm{dist}}
\newcommand{\argmin}{\operatorname{arg\,min}}
\newcommand{\argmax}{\operatorname{arg\,max}}

%% ── Header / Footer ──────────────────────────────────────────────────────
\pagestyle{fancy}
\fancyhf{}
\fancyhead[L]{\small\textit{Mathematics and Algorithms of Vector Databases}}
\fancyhead[R]{\small\thepage}
\renewcommand{\headrulewidth}{0.4pt}

%% ── Title ─────────────────────────────────────────────────────────────────
\title{
  \vspace{-1cm}
  \textbf{\LARGE Mathematics and Algorithms Behind\\[4pt]
  Vector Database Models}\\[10pt]
}


\begin{document}
\maketitle
\thispagestyle{empty}


%% ══════════════════════════════════════════════════════════════════════════
\section{Introduction and Motivation}
%% ══════════════════════════════════════════════════════════════════════════

The ``curse of dimensionality'' makes exact nearest-neighbour search
intractable in high-dimensional spaces: given $n$ points in $\R^d$ and a
query $q\in\R^d$, na\"ive linear scan requires $\calO(nd)$ time.  When
$n$ is in the hundreds of millions and $d$ ranges from 128 to 1536 (as
is common with modern neural embeddings), this cost is prohibitive.

\begin{definition}[Nearest Neighbour Search --- NNS]
Given a dataset $X = \{x_1, x_2, \ldots, x_n\} \subset \R^d$, a
distance function $\delta: \R^d \times \R^d \to \R_{\geq 0}$, and a
query point $q \in \R^d$, the \emph{nearest neighbour} of $q$ in $X$ is
\[
  \mathrm{NN}(q) = \argmin_{x \in X}\; \delta(q, x).
\]
The $k$-nearest neighbour variant ($k$-NN) returns the $k$ points with
smallest distance.
\end{definition}

\begin{definition}[$c$-Approximate Nearest Neighbour --- $c$-ANN]
For $c \geq 1$, a data structure solves $c$-ANN if, for any query $q$,
it returns a point $x^*$ such that
\[
  \delta(q, x^*) \;\leq\; c \cdot \delta\bigl(q, \mathrm{NN}(q)\bigr).
\]
\end{definition}

Vector databases implement $c$-ANN (often with $c$ very close to 1) by
combining clever indexing with optional lossy compression.  The remainder
of this document is organised as follows: Section~\ref{sec:metrics}
formalises distance metrics; Section~\ref{sec:flat} covers brute-force
search; Sections~\ref{sec:tree}--\ref{sec:hybrid} treat progressively
more sophisticated algorithms; and Section~\ref{sec:practical} discusses
practical system design.


%% ══════════════════════════════════════════════════════════════════════════
\section{Distance and Similarity Metrics}\label{sec:metrics}
%% ══════════════════════════════════════════════════════════════════════════

The choice of metric $\delta$ profoundly impacts both correctness and
efficiency.

\subsection{Euclidean ($L^2$) Distance}

\begin{definition}
For $u, v \in \R^d$,
\[
  \delta_{L^2}(u,v) = \norm{u - v}_2 = \sqrt{\sum_{i=1}^{d}(u_i - v_i)^2}.
\]
\end{definition}

Properties: a true metric (non-negativity, identity of indiscernibles,
symmetry, triangle inequality).  Squared Euclidean distance
$\delta_{L^2}^2$ is often used to avoid the square root while preserving
ranking.

\subsection{Cosine Similarity and Angular Distance}

\begin{definition}
\[
  \mathrm{sim}_{\cos}(u,v) = \frac{\inner{u}{v}}{\norm{u}_2\;\norm{v}_2}
  = \frac{\sum_{i=1}^d u_i v_i}{\sqrt{\sum_i u_i^2}\;\sqrt{\sum_i v_i^2}}.
\]
\end{definition}

The corresponding \emph{cosine distance} is $\delta_{\cos}(u,v) = 1 -
\mathrm{sim}_{\cos}(u,v)$.  Note that cosine distance is \emph{not} a
metric in the strict sense (it violates the triangle inequality), but the
\emph{angular distance}
\[
  \delta_{\mathrm{ang}}(u,v) = \frac{1}{\pi}\arccos\!\bigl(\mathrm{sim}_{\cos}(u,v)\bigr)
\]
is a valid metric on the unit sphere.

\begin{remark}
When all vectors are $L^2$-normalised, Euclidean distance and cosine
distance are monotonically related:
\[
  \norm{u-v}_2^2 = 2\bigl(1 - \mathrm{sim}_{\cos}(u,v)\bigr)
  = 2\,\delta_{\cos}(u,v).
\]
This identity is exploited in many systems to reduce cosine search to
Euclidean search after a normalisation pre-processing step.
\end{remark}

\subsection{Inner Product (Maximum Inner Product Search --- MIPS)}

\begin{definition}
\[
  \mathrm{MIPS}(q) = \argmax_{x \in X}\; \inner{q}{x}.
\]
\end{definition}

MIPS is not directly a metric search problem.  A well-known reduction
transforms MIPS into NNS by appending an extra dimension:
\[
  \tilde{x} = \bigl(x;\; \sqrt{1 - \norm{x}^2_2}\bigr), \qquad
  \tilde{q} = (q;\; 0).
\]
Then $\norm{\tilde{q} - \tilde{x}}^2_2 = \norm{q}^2_2 - 2\inner{q}{x}
+ 1$, so minimising Euclidean distance on the augmented vectors is
equivalent to maximising the inner product.

\subsection{Manhattan ($L^1$) and Chebyshev ($L^\infty$) Distances}

\begin{align}
  \delta_{L^1}(u,v) &= \sum_{i=1}^d |u_i - v_i|, \\
  \delta_{L^\infty}(u,v) &= \max_{1\le i\le d} |u_i - v_i|.
\end{align}

These are less common in embedding search but appear in specialised
domains (e.g., $L^1$ in sparse feature spaces).

\subsection{Hamming Distance}

\begin{definition}
For binary vectors $u, v \in \{0,1\}^d$,
\[
  \delta_H(u,v) = \sum_{i=1}^{d} \mathbf{1}[u_i \neq v_i]
  = \mathrm{popcount}(u \oplus v).
\]
\end{definition}

Hamming distance is critical for binary hash codes produced by LSH and
binary quantisation schemes.  Modern CPUs compute \texttt{POPCNT} in a
single instruction, making Hamming comparisons extremely fast.


%% ══════════════════════════════════════════════════════════════════════════
\section{Flat (Brute-Force) Search}\label{sec:flat}
%% ══════════════════════════════════════════════════════════════════════════

\subsection{Algorithm}

The simplest approach scans all $n$ vectors and maintains a max-heap of
size $k$.

\begin{algorithm}[H]
\caption{Flat $k$-NN Search}
\begin{algorithmic}[1]
\Require Dataset $X = \{x_1,\ldots,x_n\}\subset\R^d$, query $q\in\R^d$, integer $k$
\Ensure $k$ nearest neighbours of $q$
\State $H \gets$ empty max-heap of capacity $k$ \Comment{keyed by distance}
\For{$i = 1$ \textbf{to} $n$}
  \State $d_i \gets \delta(q, x_i)$
  \If{$|H| < k$}
    \State $H.\text{push}(x_i, d_i)$
  \ElsIf{$d_i < H.\text{top}().\text{key}$}
    \State $H.\text{pop}()$;\; $H.\text{push}(x_i, d_i)$
  \EndIf
\EndFor
\State \Return $H$
\end{algorithmic}
\end{algorithm}

\subsection{Complexity}

\begin{itemize}[nosep]
  \item \textbf{Time:} $\calO(nd + n\log k)$ per query.
  \item \textbf{Space:} $\calO(nd)$ for storage.
  \item \textbf{Preprocessing:} $\calO(1)$ (no index built).
\end{itemize}

Although asymptotically poor, flat search provides \emph{perfect recall}
and is surprisingly competitive for small datasets ($n < 10^4$), especially
with SIMD-accelerated distance computation.


%% ══════════════════════════════════════════════════════════════════════════
\section{Tree-Based Methods}\label{sec:tree}
%% ══════════════════════════════════════════════════════════════════════════

\subsection{$k$-d Trees}

\begin{definition}[$k$-d Tree~\cite{bentley1975kdtree}]
A $k$-d tree is a binary space-partitioning tree in which each internal
node stores a splitting hyperplane orthogonal to one of the $d$
coordinate axes.  At depth $\ell$, the splitting dimension is
$\ell \bmod d$.
\end{definition}

\begin{algorithm}[H]
\caption{$k$-d Tree Construction}
\begin{algorithmic}[1]
\Function{Build}{$P$, $\mathrm{depth}$}
  \If{$|P| = 0$} \Return \textsc{Nil} \EndIf
  \State $\mathrm{axis} \gets \mathrm{depth} \bmod d$
  \State Sort $P$ by coordinate $\mathrm{axis}$
  \State $m \gets \lfloor |P|/2 \rfloor$ \Comment{median index}
  \State $\mathrm{node.point} \gets P[m]$
  \State $\mathrm{node.left} \gets \Call{Build}{P[0..m{-}1],\;\mathrm{depth}+1}$
  \State $\mathrm{node.right} \gets \Call{Build}{P[m{+}1..|P|{-}1],\;\mathrm{depth}+1}$
  \State \Return $\mathrm{node}$
\EndFunction
\end{algorithmic}
\end{algorithm}

\paragraph{Query.} Nearest-neighbour search descends the tree, pruning
sub-trees whose bounding boxes are farther than the current best
candidate.

\paragraph{Complexity.}
\begin{itemize}[nosep]
  \item Build: $\calO(n\log n)$.
  \item Query (best case): $\calO(\log n)$ for low $d$.
  \item Query (worst case): $\calO(n^{1-1/d})$ --- exponential
    degradation in $d$.
\end{itemize}

Due to the curse of dimensionality, $k$-d trees are effective only for
$d \lesssim 20$.

\subsection{Ball Trees}

Ball trees use hyper-spheres instead of axis-aligned hyperplanes.

\begin{definition}[Ball Tree~\cite{omohundro1989balltree}]
Each internal node $v$ stores a centre $c_v \in \R^d$ and a radius
$r_v \geq 0$ such that all descendant points lie within the ball
$B(c_v, r_v)$.  The children partition the points into two clusters.
\end{definition}

Pruning rule: a ball $B(c_v, r_v)$ can be pruned if
$\delta(q, c_v) - r_v > d_{\text{best}}$, where $d_{\text{best}}$ is
the current best distance.

Ball trees handle general metrics (not just $L^2$) and are more
resilient to moderate dimensionality than $k$-d trees, but still
suffer at very high $d$.

\subsection{Randomised Partition Trees (Annoy)}

The \emph{Approximate Nearest Neighbours Oh Yeah} (Annoy) library~\cite{bernhardsson2015annoy} builds
a forest of randomised binary trees~\cite{dasgupta2008random}.  Each split is defined by a random
hyperplane passing through the midpoint of two randomly selected data
points.

\begin{definition}[Random Projection Split]
Given two points $a, b \in \R^d$, the splitting hyperplane has normal
$w = a - b$ and passes through the midpoint $\frac{a+b}{2}$.  A point
$x$ goes to the left child if $\inner{w}{x - \frac{a+b}{2}} \leq 0$.
\end{definition}

At query time, Annoy explores multiple trees simultaneously using a
priority queue, merging candidate sets.  Building $T$ trees on $n$
points costs $\calO(Tn\log n)$; querying has empirical complexity
$\calO(T\,d\,\log n)$ per query.


%% ══════════════════════════════════════════════════════════════════════════
\section{Locality-Sensitive Hashing (LSH)}\label{sec:lsh}
%% ══════════════════════════════════════════════════════════════════════════

\subsection{Formal Framework}

The LSH framework was introduced by Indyk and Motwani~\cite{indyk1998lsh}
and later refined by Datar et al.~\cite{datar2004e2lsh}.

\begin{definition}[$(r, cr, p_1, p_2)$-Sensitive Hash Family~\cite{indyk1998lsh}]
A family $\mathcal{H}$ of functions $h: \R^d \to \mathcal{U}$ is
$(r, cr, p_1, p_2)$-sensitive for distance $\delta$ if for any
$u, v \in \R^d$:
\begin{align}
  \delta(u,v) \leq r   &\;\Longrightarrow\; \Pr[h(u) = h(v)] \geq p_1, \\
  \delta(u,v) \geq cr  &\;\Longrightarrow\; \Pr[h(u) = h(v)] \leq p_2,
\end{align}
with $p_1 > p_2$ and $c > 1$.
\end{definition}

The \emph{LSH exponent} $\rho = \frac{\ln p_1}{\ln p_2} \in (0,1)$
governs the space--time trade-off.

\subsection{LSH for Angular Distance (SimHash / Hyperplane LSH)}

SimHash was introduced by Charikar~\cite{charikar2002simhash}.

\begin{definition}[SimHash~\cite{charikar2002simhash}]
Draw $r \sim \mathcal{N}(0, I_d)$.  Define
\[
  h_r(x) = \mathrm{sign}\bigl(\inner{r}{x}\bigr) \in \{+1, -1\}.
\]
Then $\Pr[h_r(u) = h_r(v)] = 1 - \frac{\theta(u,v)}{\pi}$, where
$\theta(u,v) = \arccos(\mathrm{sim}_{\cos}(u,v))$.
\end{definition}

By concatenating $K$ independent hash bits to form a \emph{signature} and
using $L$ independent hash tables, the query procedure becomes:

\begin{algorithm}[H]
\caption{LSH Query ($L$ tables, $K$ bits each)}
\begin{algorithmic}[1]
\Require Query $q$, hash functions $\{h_{l,1},\ldots,h_{l,K}\}_{l=1}^L$, hash tables $\{T_l\}_{l=1}^L$
\State $\mathrm{candidates} \gets \emptyset$
\For{$l = 1$ \textbf{to} $L$}
  \State $g \gets \bigl(h_{l,1}(q), \ldots, h_{l,K}(q)\bigr)$ \Comment{$K$-bit signature}
  \State $\mathrm{candidates} \gets \mathrm{candidates} \cup T_l[g]$
\EndFor
\State \Return $k\text{-NN}(q, \mathrm{candidates})$ via brute-force
\end{algorithmic}
\end{algorithm}

\subsection{LSH for Euclidean Distance (E2LSH)}

\begin{definition}[E2LSH Hash~\cite{datar2004e2lsh}]
Let $a \sim \mathcal{N}(0, I_d)$ and $b \sim \mathrm{Uniform}[0, w)$
for bucket width $w > 0$.  Then
\[
  h_{a,b}(x) = \left\lfloor \frac{\inner{a}{x} + b}{w} \right\rfloor.
\]
\end{definition}

The collision probability depends on $\delta_{L^2}(u,v)/w$ and has a
closed-form expression involving the Gaussian CDF.

\subsection{Complexity}

\begin{theorem}
For an $(r, cr, p_1, p_2)$-sensitive family with $\rho = \frac{\ln
p_1}{\ln p_2}$, there exists an LSH data structure solving $c$-ANN with:
\begin{itemize}[nosep]
  \item Space: $\calO(n^{1+\rho} + nd)$.
  \item Query time: $\calO(n^{\rho}\,d)$.
  \item Preprocessing: $\calO(n^{1+\rho}\,d)$.
\end{itemize}
\end{theorem}

For SimHash on angular distance with approximation ratio $c$,
$\rho = 1/c$.

\subsection{Multi-Probe LSH}

Standard LSH requires many tables.  Multi-probe LSH~\cite{lv2007multiprobe} reduces this by
probing \emph{neighbouring} buckets---those obtained by flipping one or
more hash bits---using a perturbation scheme ranked by expected distance.
This can reduce $L$ by an order of magnitude while maintaining recall.


%% ══════════════════════════════════════════════════════════════════════════
\section{Quantisation Methods}\label{sec:quant}
%% ══════════════════════════════════════════════════════════════════════════

\subsection{Scalar Quantisation (SQ)}

\begin{definition}[Scalar Quantisation]
Each floating-point coordinate $x_i \in \R$ is independently mapped to
one of $2^b$ levels.  For uniform quantisation with $b$ bits:
\[
  \hat{x}_i = \mathrm{round}\!\left(
    \frac{x_i - x_{\min}}{x_{\max} - x_{\min}} \cdot (2^b - 1)
  \right).
\]
\end{definition}

Common choices are $b \in \{4, 8\}$.  With $b=8$, storage shrinks from
$4d$ bytes (float32) to $d$ bytes per vector.  Distance computation uses
look-up tables or integer SIMD.

\subsection{Product Quantisation (PQ)}\label{subsec:pq}

Product quantisation, introduced by J\'egou et al.~\cite{jegou2011pq}, is the
workhorse compression scheme for large-scale vector search.

\begin{definition}[Product Quantisation~\cite{jegou2011pq}]
Decompose $\R^d$ into $M$ disjoint sub-spaces of dimension $d/M$ each:
\[
  x = \bigl(x^{(1)}, x^{(2)}, \ldots, x^{(M)}\bigr), \qquad
  x^{(m)} \in \R^{d/M}.
\]
In each sub-space, learn a codebook $C^{(m)} = \{c^{(m)}_1, \ldots,
c^{(m)}_{K_s}\}$ via $k$-means (typically $K_s = 256$).  A vector $x$
is encoded as
\[
  \mathrm{PQ}(x) = \bigl(q^{(1)}(x), \ldots, q^{(M)}(x)\bigr),
  \quad q^{(m)}(x) = \argmin_{j}\; \norm{x^{(m)} - c^{(m)}_j}_2.
\]
Each code $q^{(m)}(x) \in \{1,\ldots,K_s\}$ requires $\lceil\log_2
K_s\rceil = 8$ bits for $K_s=256$, so the total code is $M$ bytes.
\end{definition}

\paragraph{Asymmetric Distance Computation (ADC).}
At query time, we compute exact sub-space distances between $q$ and each
centroid, then approximate the full distance:
\[
  \hat{\delta}(q, x)^2 = \sum_{m=1}^{M}
  \norm{q^{(m)} - c^{(m)}_{q^{(m)}(x)}}_2^2.
\]
The $M \times K_s$ distance table is precomputed in $\calO(MK_s \cdot
d/M) = \calO(K_s d)$ time.  Then each vector comparison costs only
$\calO(M)$ look-ups.

\begin{algorithm}[H]
\caption{PQ Search with ADC}
\begin{algorithmic}[1]
\Require Codebooks $\{C^{(m)}\}_{m=1}^M$, codes $\{\mathrm{PQ}(x_i)\}_{i=1}^n$, query $q$
\For{$m = 1$ \textbf{to} $M$} \Comment{Build distance table}
  \For{$j = 1$ \textbf{to} $K_s$}
    \State $D[m][j] \gets \norm{q^{(m)} - c^{(m)}_j}_2^2$
  \EndFor
\EndFor
\For{$i = 1$ \textbf{to} $n$} \Comment{Scan all codes}
  \State $\hat{d}_i \gets \sum_{m=1}^M D\bigl[m\bigr]\bigl[\mathrm{PQ}(x_i)_m\bigr]$
\EndFor
\State \Return top-$k$ by smallest $\hat{d}_i$
\end{algorithmic}
\end{algorithm}

\subsection{Optimised Product Quantisation (OPQ)~\cite{ge2013opq}}

Standard PQ assumes the sub-spaces are axis-aligned.  OPQ~\cite{ge2013opq} learns a
rotation matrix $R \in \R^{d \times d}$ such that the rotated data
$Rx$ has minimised quantisation error when partitioned into sub-spaces:
\[
  R^* = \argmin_{R \in O(d)} \sum_{i=1}^n
  \norm{Rx_i - \widetilde{Rx_i}}_2^2,
\]
where $\widetilde{Rx_i}$ is the PQ reconstruction of $Rx_i$.  This is
solved by alternating between codebook updates (standard $k$-means) and
rotation updates (Procrustes).

\subsection{Additive Quantisation (AQ) and Residual Quantisation (RQ)}

\begin{definition}[Additive Quantisation~\cite{babenko2014aq}]
Represent each vector as a sum of codewords from $M$ codebooks (which
now span the \emph{full} $d$-dimensional space):
\[
  \hat{x} = \sum_{m=1}^{M} c^{(m)}_{j_m}, \qquad c^{(m)}_{j_m} \in \R^d.
\]
The codes $(j_1, \ldots, j_M)$ are chosen jointly to minimise
$\norm{x - \hat{x}}_2^2$.
\end{definition}

This is more expressive than PQ (since codebooks interact across
dimensions) but encoding is NP-hard; beam search is used in practice.

\begin{definition}[Residual Quantisation~\cite{chen2010rq}]
A greedy special case of AQ: quantise $x$ with the first codebook to
get residual $r_1 = x - c^{(1)}_{j_1}$, then quantise $r_1$ with the
second codebook, and so on:
\[
  r_0 = x, \qquad j_m = \argmin_j \norm{r_{m-1} - c^{(m)}_j}, \qquad
  r_m = r_{m-1} - c^{(m)}_{j_m}.
\]
\end{definition}

RQ often achieves better distortion than PQ at the same code length,
especially when combined with a learnt rotation (RQ with OPQ).

\subsection{Binary Quantisation}

The simplest quantisation: $\hat{x}_i = \mathrm{sign}(x_i)$.  Reduces
$d$-dimensional float32 vectors to $d$-bit strings.  Distances are
approximated via Hamming distance, enabling extremely fast comparison
($\sim$1 clock cycle per 64 bits via \texttt{POPCNT}).  However,
the information loss is severe; binary quantisation is typically used
only as a \emph{re-ranking filter}~\cite{aguerrebere2024similarity}: retrieve a large candidate set via
Hamming, then re-rank with exact distances.


%% ══════════════════════════════════════════════════════════════════════════
\section{Inverted File Index (IVF)}\label{sec:ivf}
%% ══════════════════════════════════════════════════════════════════════════

\subsection{Core Idea}

Partition the dataset into $C$ Voronoi cells via $k$-means clustering
with centroids $\mu_1, \ldots, \mu_C$.  Each vector $x$ is assigned to
its closest centroid.  At query time, only the $n_{\text{probe}}$ closest
cells are searched.

\begin{definition}[IVF Assignment]
\[
  \mathrm{cell}(x) = \argmin_{c \in \{1,\ldots,C\}}\; \delta(x, \mu_c).
\]
\end{definition}

\begin{algorithm}[H]
\caption{IVF Search}
\begin{algorithmic}[1]
\Require Centroids $\{\mu_c\}_{c=1}^C$, inverted lists $\{L_c\}_{c=1}^C$, query $q$, $n_{\mathrm{probe}}$
\State $S \gets n_{\mathrm{probe}}$ closest centroids to $q$
\State $\mathrm{candidates} \gets \bigcup_{c \in S} L_c$
\State \Return $k\text{-NN}(q, \mathrm{candidates})$
\end{algorithmic}
\end{algorithm}

\subsection{IVF-PQ: Combining IVF with Product Quantisation}

Store PQ codes instead of full vectors in each inverted list.  The
residual $r_x = x - \mu_{\mathrm{cell}(x)}$ is quantised, improving
distortion.  This is the backbone of the \textsc{Faiss} library's~\cite{johnson2019faiss} most
widely-used index (\texttt{IVF$C$,PQ$M$}).

\paragraph{Complexity.}
\begin{itemize}[nosep]
  \item Index build: $\calO(n d C_{\text{iter}})$ for $k$-means +
    $\calO(nd)$ for PQ encoding.
  \item Query: $\calO(Cd + n_{\text{probe}} \cdot
    \frac{n}{C} \cdot M)$ (centroid comparison + PQ scan).
  \item Space: $\calO(Cd + nM)$.
\end{itemize}

By varying $n_{\text{probe}}$, one tunes the recall/speed trade-off
continuously.

\subsection{IVF-Flat}

Same as IVF but stores full vectors.  Higher memory but exact distances
within each probed cell.  Useful when $n$ is moderate ($10^5$--$10^6$).


%% ══════════════════════════════════════════════════════════════════════════
\section{Graph-Based Methods}\label{sec:graph}
%% ══════════════════════════════════════════════════════════════════════════

Graph-based indices construct a proximity graph $G = (V, E)$ where
$V = X$ and edges connect nearby points.  Queries are answered by a
greedy walk from an entry point, following edges that decrease the
distance to $q$.

\subsection{Navigable Small World (NSW)}

The NSW graph construction was proposed by Malkov et al.~\cite{malkov2014nsw}.

\begin{definition}[NSW Graph~\cite{malkov2014nsw}]
An NSW graph is built by inserting points one at a time.  When inserting
point $x$, a greedy search from a random entry point finds the current
approximate nearest neighbours; edges are added between $x$ and these
neighbours.
\end{definition}

The graph exhibits small-world properties: short-range edges ensure local
accuracy while long-range edges (created by early insertions) provide
logarithmic routing.

\subsection{Hierarchical Navigable Small World (HNSW)}\label{subsec:hnsw}

HNSW, proposed by Malkov and Yashunin~\cite{malkov2018hnsw}, layers multiple NSW graphs
at different resolutions.

\begin{definition}[HNSW Structure]
An HNSW index has $L_{\max}$ layers, $\ell = 0, 1, \ldots, L_{\max}$.
Each point is assigned a random maximum layer
\[
  l(x) = \lfloor -\ln(\mathrm{Uniform}(0,1)) \cdot m_L \rfloor,
\]
where $m_L = 1/\ln(M)$ and $M$ is a construction parameter.
Point $x$ appears in all layers $0, 1, \ldots, l(x)$.  Layer 0
(the bottom) contains all points; upper layers are exponentially
sparser.
\end{definition}

\begin{algorithm}[H]
\caption{HNSW Search}
\begin{algorithmic}[1]
\Require HNSW graph $G$, query $q$, $k$, $\mathrm{ef}$ (beam width)
\State $\mathrm{ep} \gets$ entry point (top layer node)
\For{$\ell = L_{\max}$ \textbf{downto} $1$} \Comment{Greedy descent}
  \State $\mathrm{ep} \gets \text{GreedySearch}(q, \mathrm{ep}, 1, \ell)$
\EndFor
\State $W \gets \text{BeamSearch}(q, \mathrm{ep}, \mathrm{ef}, 0)$
  \Comment{Layer 0: beam search}
\State \Return closest $k$ elements from $W$
\end{algorithmic}
\end{algorithm}

\begin{algorithm}[H]
\caption{HNSW Insertion}
\begin{algorithmic}[1]
\Require New point $x$, HNSW graph $G$, parameters $M$, $M_{\max}$, $\mathrm{efConstruction}$
\State $l \gets \lfloor -\ln(\mathrm{Uniform}(0,1)) / \ln(M) \rfloor$
\State $\mathrm{ep} \gets$ current entry point
\For{$\ell = L_{\max}$ \textbf{downto} $l+1$}
  \State $\mathrm{ep} \gets \text{GreedySearch}(x, \mathrm{ep}, 1, \ell)$
\EndFor
\For{$\ell = \min(l, L_{\max})$ \textbf{downto} $0$}
  \State $W \gets \text{BeamSearch}(x, \mathrm{ep}, \mathrm{efConstruction}, \ell)$
  \State Select $M$ neighbours from $W$ (heuristic or simple nearest)
  \State Add bidirectional edges; prune if degree exceeds $M_{\max}$
  \State $\mathrm{ep} \gets$ closest element in $W$
\EndFor
\If{$l > L_{\max}$}
  \State Set $x$ as new entry point; $L_{\max} \gets l$
\EndIf
\end{algorithmic}
\end{algorithm}

\paragraph{Neighbour Selection Heuristic.}
Rather than selecting the $M$ closest points, HNSW uses a heuristic
that promotes \emph{diverse} connections.  Candidates are processed in
order of distance; a candidate is added only if it is closer to $x$ than
to any already-selected neighbour.  This ensures edges span gaps in the
data distribution, improving connectivity and recall.

\paragraph{Complexity.}
\begin{itemize}[nosep]
  \item Construction: $\calO(n \cdot \log n \cdot M \cdot d)$ amortised.
  \item Query: $\calO(\mathrm{ef} \cdot M \cdot d \cdot \log n)$
    empirically; $\mathrm{ef}$ controls the recall/speed trade-off.
  \item Space: $\calO(nMd + nM\log n)$ for edges and vectors.
\end{itemize}

HNSW achieves state-of-the-art recall at high throughput and is the
default index in libraries such as \textsc{hnswlib}, \textsc{Faiss}~\cite{douze2024faiss},
\textsc{Milvus}, \textsc{Weaviate}, and \textsc{Qdrant}.


%% ══════════════════════════════════════════════════════════════════════════
\section{Vamana and DiskANN}\label{sec:diskann}
%% ══════════════════════════════════════════════════════════════════════════

\subsection{Motivation}

HNSW requires the entire graph and all vectors in RAM.  For billion-scale
datasets this can require terabytes of memory.  \emph{DiskANN}~\cite{jayaram2019diskann} stores the graph on SSD while keeping a
compressed representation in memory.

\subsection{Vamana Graph Construction}

Vamana builds a single-layer directed graph with bounded out-degree $R$.

\begin{algorithm}[H]
\caption{Vamana Index Construction}
\begin{algorithmic}[1]
\Require Dataset $X$, max degree $R$, beam width $L$, parameter $\alpha \geq 1$
\State Initialise $G$ with random $R$-regular graph
\State $s \gets$ medoid of $X$ \Comment{entry point}
\For{each $x \in X$ (random order, two passes)}
  \State $\mathcal{V} \gets \text{GreedySearch}(x, s, L)$ \Comment{visited list}
  \State $N \gets \text{RobustPrune}(x, \mathcal{V}, \alpha, R)$
  \State Set out-neighbours of $x$ to $N$
  \For{each $p \in N$}
    \If{$|\text{out}(p)| > R$}
      \State $\text{out}(p) \gets \text{RobustPrune}(p, \text{out}(p) \cup \{x\}, \alpha, R)$
    \Else
      \State $\text{out}(p) \gets \text{out}(p) \cup \{x\}$
    \EndIf
  \EndFor
\EndFor
\end{algorithmic}
\end{algorithm}

\paragraph{Robust Prune.}
The \textsc{RobustPrune} procedure generalises HNSW's heuristic
neighbour selection.  Given a candidate set $\mathcal{V}$ and a distance
threshold factor $\alpha$, it iteratively selects the closest
unprocessed candidate $p^*$ and removes all candidates $p$ satisfying
$\alpha \cdot \delta(p^*, p) \leq \delta(x, p)$.  The parameter
$\alpha > 1$ encourages long-range edges, improving navigability.

\subsection{DiskANN Query}

\begin{enumerate}[nosep]
  \item Maintain a PQ-compressed copy of all vectors in memory.
  \item Greedy search on the graph using PQ distances (fast, approximate).
  \item For the most promising candidates, issue SSD reads to fetch full
    vectors and re-rank with exact distances.
\end{enumerate}

With NVMe SSDs providing $\sim$500K random 4\,KB reads/sec, DiskANN
achieves sub-millisecond query latency on billion-point datasets using
only 64\,GB of RAM + SSD.


%% ══════════════════════════════════════════════════════════════════════════
\section{ScaNN: Score-Aware Quantisation}\label{sec:scann}
%% ══════════════════════════════════════════════════════════════════════════

Google's \emph{ScaNN} (Scalable Nearest Neighbours)~\cite{guo2020scann}
introduces \emph{anisotropic vector quantisation}, which accounts for the
direction of the query during codebook learning.

\subsection{Anisotropic Loss}

Standard PQ minimises the \emph{reconstruction error}:
$\sum_i \norm{x_i - \hat{x}_i}^2$.  ScaNN instead minimises a loss
that weights errors along the query direction more heavily:
\[
  \mathcal{L}_{\mathrm{ScaNN}} = \sum_{i=1}^n \left[
    \bigl(\inner{x_i}{\hat{x}_i} - \norm{x_i}^2\bigr)^2
    + \lambda \norm{x_i - \hat{x}_i}_{\perp}^2
  \right],
\]
where $\norm{\cdot}_{\perp}$ denotes the component orthogonal to $x_i$
and $\lambda < 1$ down-weights it.  This yields quantised codes that
better preserve inner products (relevant for MIPS).

\subsection{Architecture}

ScaNN combines three stages:
\begin{enumerate}[nosep]
  \item \textbf{Partitioning}: $k$-means (IVF-style) coarse quantiser.
  \item \textbf{Scoring}: Anisotropic PQ for fast approximate scoring.
  \item \textbf{Re-ranking}: Exact distance computation on the top
    candidates.
\end{enumerate}


%% ══════════════════════════════════════════════════════════════════════════
\section{Hybrid and Filtered Search}\label{sec:hybrid}
%% ══════════════════════════════════════════════════════════════════════════

\subsection{Pre-Filtering vs.\ Post-Filtering}

In many applications, queries carry metadata predicates (e.g., ``find
similar images \emph{uploaded after 2023}'').  Two strategies:

\begin{itemize}[nosep]
  \item \textbf{Post-filter}: Run ANN, then discard results failing the
    predicate.  Fast but recall degrades when the predicate is selective.
  \item \textbf{Pre-filter}: Identify the subset satisfying the
    predicate, then search within it.  Exact but may invalidate index
    structure.
\end{itemize}

\subsection{Filtered Graph Search}

Modern systems (e.g., Qdrant, Weaviate) integrate filtering into the
graph traversal itself.  During greedy search on HNSW, a visited
neighbour that fails the predicate is \emph{not added to the result set}
but \emph{is still used for navigation} (its neighbours are explored).
This preserves the connectivity of the graph while enforcing the filter.

\subsection{Hybrid IVF}

Partition vectors into cells \emph{and} maintain per-cell metadata
bitmaps.  At query time, intersect the bitmap with the predicate before
scanning each cell.

\subsection{Multi-Vector Search and Late Interaction}

Models like ColBERT~\cite{khattab2020colbert} represent a document as a \emph{set} of token-level
vectors.  The relevance score uses \emph{MaxSim}:
\[
  \mathrm{score}(q, d) = \sum_{i=1}^{|q|} \max_{j=1}^{|d|}
  \inner{q_i}{d_j}.
\]
Vector databases supporting multi-vector search index individual token
vectors and aggregate scores at query time.


%% ══════════════════════════════════════════════════════════════════════════
\section{Dimensionality Reduction Techniques}\label{sec:dimred}
%% ══════════════════════════════════════════════════════════════════════════

Reducing $d$ before indexing accelerates all downstream operations.

\subsection{PCA (Principal Component Analysis)}

Project onto the top-$k$ eigenvectors of the covariance matrix:
\[
  \hat{x} = W_k^\top (x - \bar{x}), \qquad W_k \in \R^{d \times k},
\]
where $W_k$ consists of the $k$ eigenvectors with largest eigenvalues
of $\frac{1}{n}\sum_i (x_i - \bar{x})(x_i - \bar{x})^\top$.

PCA preserves the maximum variance and is optimal for Euclidean
reconstruction error under a linear map.

\subsection{Random Projection (Johnson--Lindenstrauss)}

\begin{theorem}[Johnson--Lindenstrauss Lemma~\cite{johnson1984extensions}]
For any $\epsilon \in (0,1)$ and any set of $n$ points in $\R^d$, there
exists a linear map $f: \R^d \to \R^k$ with
$k = \calO(\epsilon^{-2} \log n)$ such that for all pairs $(u, v)$:
\[
  (1 - \epsilon)\norm{u-v}_2^2 \leq \norm{f(u) - f(v)}_2^2
  \leq (1+\epsilon)\norm{u-v}_2^2.
\]
Moreover, $f$ can be a random matrix with i.i.d.\ $\mathcal{N}(0,1/k)$
entries.
\end{theorem}

We now give a complete proof, following the standard approach via
sub-Gaussian concentration.  The argument proceeds in three stages:
(1)~analyse a single projected squared norm, (2)~obtain a two-sided
tail bound, and (3)~take a union bound over all pairs.

\begin{proof}
\textbf{Setup.}\;
Let $A \in \R^{k \times d}$ be a random matrix whose entries are
i.i.d.\ $A_{ij} \sim \mathcal{N}(0, 1/k)$.  Define the projection
$f(x) = Ax$ for $x \in \R^d$.  By linearity it suffices to show that
for any \emph{fixed} unit vector $u \in \R^d$ (we will later set
$u = (v_1 - v_2)/\norm{v_1 - v_2}_2$ for each pair),
\[
  \Pr\!\Bigl[\bigl|\norm{Au}_2^2 - 1\bigr| > \epsilon\Bigr]
  \;\leq\; 2\exp\!\left(-\frac{k(\epsilon^2 - \epsilon^3)}{4}\right).
  \tag{$\star$}
\]

\medskip
\textbf{Step 1: Distribution of the projected squared norm.}\;
Write $Y = Au \in \R^k$.  Because $u$ is a fixed unit vector and the
rows of $A$ are independent,
\[
  Y_i = \sum_{j=1}^{d} A_{ij}\, u_j, \qquad i = 1, \ldots, k.
\]
Each $Y_i$ is a linear combination of independent Gaussians with
variance $\sum_j u_j^2 / k = 1/k$.  Hence $Y_i \sim
\mathcal{N}(0, 1/k)$ and the $Y_i$ are mutually independent.
Therefore
\[
  \norm{Au}_2^2 = \sum_{i=1}^{k} Y_i^2
  = \frac{1}{k}\sum_{i=1}^{k} Z_i^2,
\]
where $Z_i = \sqrt{k}\, Y_i \sim \mathcal{N}(0,1)$ are i.i.d.\ standard
normals, so $\sum_i Z_i^2 \sim \chi^2_k$.  Define
$W = \norm{Au}_2^2 = \frac{1}{k}\chi^2_k$.  We have $\mathbb{E}[W] = 1$.

\medskip
\textbf{Step 2: Sub-exponential tail bound via the moment-generating
function.}\;
For the upper tail, let $t > 0$.  By Markov's inequality applied to the
MGF:
\[
  \Pr[W > 1 + \epsilon]
  = \Pr\!\bigl[e^{tW} > e^{t(1+\epsilon)}\bigr]
  \leq \frac{\mathbb{E}[e^{tW}]}{e^{t(1+\epsilon)}}.
\]
Since $W = \frac{1}{k}\sum_i Z_i^2$ and the $Z_i$ are independent,
\[
  \mathbb{E}[e^{tW}]
  = \prod_{i=1}^{k} \mathbb{E}\!\left[e^{t Z_i^2 / k}\right].
\]
For $Z \sim \mathcal{N}(0,1)$ and $s < 1/2$,
\[
  \mathbb{E}[e^{sZ^2}] = \frac{1}{\sqrt{1-2s}}.
\]
Setting $s = t/k$ (requiring $t < k/2$):
\[
  \mathbb{E}[e^{tW}] = (1 - 2t/k)^{-k/2}.
\]
Hence
\[
  \Pr[W > 1+\epsilon]
  \leq (1 - 2t/k)^{-k/2}\, e^{-t(1+\epsilon)}.
\]
Optimising over $t$: set $t = \frac{k\epsilon}{2(1+\epsilon)}$ (which
satisfies $t < k/2$).  Substituting:
\begin{align}
  \ln\Pr[W > 1+\epsilon]
  &\leq \frac{k}{2}\left[
    \epsilon - \ln(1+\epsilon)
  \right] \notag\\
  &\leq \frac{k}{2}\left[
    \epsilon - \epsilon + \tfrac{\epsilon^2}{2} - \tfrac{\epsilon^3}{3} + \cdots
  \right] \cdot (-1) \notag\\
  &= -\frac{k}{2}\left[
    \ln(1+\epsilon) - \epsilon
  \right]. \label{eq:jl-upper}
\end{align}
Using the elementary inequality $\ln(1+\epsilon) \geq \epsilon -
\epsilon^2/2$ for $\epsilon \in (0,1)$, we get
\[
  \Pr[W > 1+\epsilon] \leq \exp\!\left(-\frac{k\epsilon^2}{4}\right).
\]

For a tighter bound that will be useful for the final constant, we use
$\ln(1+\epsilon) \geq \epsilon - \epsilon^2/2 + \epsilon^3/3$ (valid
for $0 < \epsilon < 1$), yielding the refinement:
\[
  \Pr[W > 1+\epsilon] \leq \exp\!\left(-\frac{k(\epsilon^2/2 - \epsilon^3/3)}{2}\right)
  = \exp\!\left(-\frac{k(\epsilon^2 - 2\epsilon^3/3)}{4}\right).
\]

\medskip
\textbf{Lower tail.}\;
For the lower tail $\Pr[W < 1 - \epsilon]$, an analogous calculation
with $t < 0$ (or equivalently, using $\mathbb{E}[e^{-tW}]$) gives
\begin{align}
  \Pr[W < 1-\epsilon]
  &\leq \exp\!\left(-\frac{k}{2}\bigl[-\epsilon - \ln(1-\epsilon)\bigr]\right) \notag\\
  &\leq \exp\!\left(-\frac{k(\epsilon^2 - \epsilon^3)}{4}\right),
  \label{eq:jl-lower}
\end{align}
where we used $-\ln(1-\epsilon) \geq \epsilon + \epsilon^2/2$ and retained
a cubic correction.  (The lower tail bound is slightly weaker because
$(1-\epsilon)$ is farther from the mode than $(1+\epsilon)$ in relative
terms.)

Combining the two tails via a union bound:
\[
  \Pr\!\bigl[\bigl|\norm{Au}_2^2 - 1\bigr| > \epsilon\bigr]
  \leq 2\exp\!\left(-\frac{k(\epsilon^2 - \epsilon^3)}{4}\right).
\]
This establishes $(\star)$.

\medskip
\textbf{Step 3: Union bound over all pairs.}\;
There are $\binom{n}{2} < n^2/2$ pairs among $n$ points.  For a fixed
pair $(v_1, v_2)$, setting $u = (v_1-v_2)/\norm{v_1-v_2}_2$:
\[
  \norm{f(v_1) - f(v_2)}_2^2 = \norm{v_1 - v_2}_2^2 \cdot \norm{Au}_2^2,
\]
so the event $\bigl|\norm{Au}_2^2 - 1\bigr| \leq \epsilon$ is
equivalent to
\[
  (1-\epsilon)\norm{v_1-v_2}_2^2
  \leq \norm{f(v_1)-f(v_2)}_2^2
  \leq (1+\epsilon)\norm{v_1-v_2}_2^2.
\]
By the union bound, the probability that \emph{any} pair violates this
is at most
\[
  \binom{n}{2} \cdot 2\exp\!\left(-\frac{k(\epsilon^2 - \epsilon^3)}{4}\right)
  < n^2 \exp\!\left(-\frac{k(\epsilon^2 - \epsilon^3)}{4}\right).
\]
Setting this less than a target failure probability $\delta > 0$:
\[
  n^2 \exp\!\left(-\frac{k(\epsilon^2 - \epsilon^3)}{4}\right) < \delta
  \quad\Longleftrightarrow\quad
  k > \frac{4(\,2\ln n - \ln \delta\,)}{\epsilon^2 - \epsilon^3}.
\]
For $\epsilon \in (0, 1/2]$ we have $\epsilon^2 - \epsilon^3 \geq
\epsilon^2/2$, so
\[
  k = \left\lceil \frac{8(\,2\ln n - \ln\delta\,)}{\epsilon^2} \right\rceil
  = \calO\!\left(\frac{\log n}{\epsilon^2}\right)
\]
suffices (with $\delta$ taken as any fixed constant, e.g.\ $\delta = 1/n$).

\medskip
\textbf{Conclusion.}\;
With probability at least $1 - \delta$, the random linear map
$f(x) = Ax$ with $A \in \R^{k \times d}$, \;$A_{ij} \sim
\mathcal{N}(0,1/k)$, and
$k = \calO(\epsilon^{-2}\log n)$
satisfies
\[
  (1-\epsilon)\norm{u-v}_2^2 \leq \norm{f(u)-f(v)}_2^2
  \leq (1+\epsilon)\norm{u-v}_2^2
\]
simultaneously for all $\binom{n}{2}$ pairs.  In particular, the
existence of such a map is guaranteed (probabilistic method), and
moreover a random Gaussian matrix achieves it with high probability.
\end{proof}

\begin{remark}[Tightness]
The dependence $k = \Theta(\epsilon^{-2}\log n)$ is known to be
\emph{optimal}: Alon~\cite{alon2003jl} showed that for certain point sets, any
linear map into $\R^k$ preserving all pairwise distances to within
$(1 \pm \epsilon)$ requires $k = \Omega(\epsilon^{-2}\log n)$.
\end{remark}

\begin{remark}[Generalisations]
The Gaussian distribution is not essential.  The proof uses only
two properties of the entries of $A$: (i)~zero mean and variance $1/k$,
and (ii)~sub-Gaussian tails.  Consequently, the lemma holds with the
same $k$ for:
\begin{itemize}[nosep]
  \item \textbf{Rademacher entries:} $A_{ij} = \pm 1/\sqrt{k}$ with
    equal probability~\cite{achlioptas2003random}.
  \item \textbf{Sparse projections:} $A_{ij} \in
    \{-\sqrt{3/k},\; 0,\; \sqrt{3/k}\}$ with probabilities
    $\{1/6,\; 2/3,\; 1/6\}$~\cite{achlioptas2003random}, yielding $\sim$3$\times$
    speedup in projection.
  \item \textbf{Sub-sampled randomised Hadamard transform (SRHT):}
    achieves $\calO(d \log k)$ projection time instead of
    $\calO(dk)$~\cite{ailon2009fast}.
\end{itemize}
\end{remark}

\subsection{Matryoshka Representation Learning (MRL)}

A recent training technique~\cite{kusupati2022matryoshka} that produces embeddings where truncating
to the first $k < d$ dimensions yields a valid (though coarser)
embedding.  This enables adaptive precision: coarse search in low
dimensions, then re-rank in full dimensions.


%% ══════════════════════════════════════════════════════════════════════════
\section{Practical System Design Considerations}\label{sec:practical}
%% ══════════════════════════════════════════════════════════════════════════

\subsection{Index Selection Guidelines}

\begin{table}[H]
\centering
\small
\begin{tabular}{@{}lllll@{}}
\toprule
\textbf{Method} & \textbf{Build Time} & \textbf{Query Time} & \textbf{Memory} & \textbf{Best For} \\
\midrule
Flat          & $\calO(1)$    & $\calO(nd)$      & $\calO(nd)$ & $n < 10^4$ \\
IVF-Flat      & $\calO(nd)$   & $\calO(\frac{n}{C}d)$ & $\calO(nd)$ & $10^4$--$10^6$ \\
IVF-PQ        & $\calO(nd)$   & $\calO(\frac{n}{C}M)$ & $\calO(nM)$ & $10^6$--$10^9$, RAM-limited \\
HNSW          & $\calO(n\log n\,d)$ & $\calO(\mathrm{ef}\cdot d\log n)$ & $\calO(nMd)$ & $10^5$--$10^8$, latency-critical \\
DiskANN       & $\calO(n\log n\,d)$ & $\calO(d\log n)$ & $\calO(nM)$ + SSD & $>10^8$, limited RAM \\
LSH           & $\calO(n^{1+\rho}d)$ & $\calO(n^\rho d)$ & $\calO(n^{1+\rho})$ & Theoretical / streaming \\
\bottomrule
\end{tabular}
\caption{Complexity comparison of major index types ($C$ = IVF clusters, $M$ = PQ sub-spaces or HNSW degree, $\mathrm{ef}$ = HNSW beam width).}
\end{table}

\subsection{Recall--Throughput Trade-off}

Every ANN index exposes one or more ``knobs'' controlling the
recall--throughput trade-off:
\begin{itemize}[nosep]
  \item \textbf{IVF}: $n_{\text{probe}}$ (number of cells probed).
  \item \textbf{HNSW}: $\mathrm{ef}$ (beam width at query time).
  \item \textbf{PQ}: $M$ (number of sub-quantisers) and $K_s$ (codebook
    size).
  \item \textbf{DiskANN}: $L$ (search list size) and number of SSD
    reads.
\end{itemize}

\subsection{Sharding and Replication}

For datasets exceeding single-machine capacity, vector databases
partition data across shards (each shard is an independent index) and
replicate shards for fault tolerance.  Query routing identifies relevant
shards, issues parallel sub-queries, and merges top-$k$ results.

\subsection{Streaming / Real-Time Updates}

Graph-based indices (HNSW, Vamana) support incremental insertion
natively.  IVF indices require periodic re-clustering as the data
distribution drifts.  A common pattern is a \emph{two-tier} architecture:
a mutable in-memory buffer (e.g., HNSW) for recent writes, periodically
merged into an immutable on-disk index (e.g., IVF-PQ).


%% ══════════════════════════════════════════════════════════════════════════
\section{Evaluation Metrics}\label{sec:eval}
%% ══════════════════════════════════════════════════════════════════════════

\begin{definition}[Recall@$k$]
Given ground-truth $k$-NN set $G_k(q)$ and retrieved set $R_k(q)$:
\[
  \mathrm{Recall@}k = \frac{|R_k(q) \cap G_k(q)|}{k}.
\]
\end{definition}

\begin{definition}[Queries Per Second (QPS)]
Throughput metric; typically measured single-threaded or at a fixed
parallelism level.
\end{definition}

The standard benchmark is \textsc{ANN-Benchmarks}~\cite{aumuller2020annbenchmarks}, which evaluates indices on the Pareto frontier of recall vs.\ QPS
across standardised datasets (SIFT1M, GloVe, GIST, etc.).


%% ══════════════════════════════════════════════════════════════════════════
\section{Summary and Outlook}
%% ══════════════════════════════════════════════════════════════════════════

We have presented the mathematical and algorithmic foundations of vector
database systems, progressing from exact brute-force search through
tree-based methods, locality-sensitive hashing, quantisation techniques,
inverted file indices, and state-of-the-art graph-based approaches.

Key themes include the fundamental tension between recall and efficiency;
the power of quantisation to compress vectors with controlled distortion;
and the remarkable effectiveness of navigable small-world graphs for
high-dimensional search.

Active research frontiers include learned indices (using neural networks
to predict partition boundaries), hardware-aware optimisation (GPU,
FPGA, CXL-memory), streaming and mutable indices for real-time
applications, and tighter integration of vector search with relational
query processing in hybrid databases.


%% ══════════════════════════════════════════════════════════════════════════
\begin{thebibliography}{99}
%% ══════════════════════════════════════════════════════════════════════════

\bibitem{johnson1984extensions}
W.~B. Johnson and J.~Lindenstrauss,
``Extensions of Lipschitz mappings into a Hilbert space,''
\emph{Contemporary Mathematics}, vol.~26, pp.~189--206, 1984.

\bibitem{indyk1998lsh}
P.~Indyk and R.~Motwani,
``Approximate nearest neighbors: Towards removing the curse of dimensionality,''
in \emph{Proc.\ 30th ACM Symposium on Theory of Computing (STOC)},
pp.~604--613, 1998.

\bibitem{charikar2002simhash}
M.~S. Charikar,
``Similarity estimation techniques from rounding algorithms,''
in \emph{Proc.\ 34th ACM Symposium on Theory of Computing (STOC)},
pp.~380--388, 2002.

\bibitem{achlioptas2003random}
D.~Achlioptas,
``Database-friendly random projections: Johnson--Lindenstrauss with binary coins,''
\emph{Journal of Computer and System Sciences}, vol.~66, no.~4,
pp.~671--687, 2003.

\bibitem{alon2003jl}
N.~Alon,
``Problems and results in extremal combinatorics---I,''
\emph{Discrete Mathematics}, vol.~273, no.~1--3, pp.~31--53, 2003.
(Lower bound for the Johnson--Lindenstrauss lemma.)

\bibitem{datar2004e2lsh}
M.~Datar, N.~Immorlica, P.~Indyk, and V.~S. Mirrokni,
``Locality-sensitive hashing scheme based on $p$-stable distributions,''
in \emph{Proc.\ 20th ACM Symposium on Computational Geometry (SoCG)},
pp.~253--262, 2004.

\bibitem{lv2007multiprobe}
Q.~Lv, W.~Josephson, Z.~Wang, M.~Charikar, and K.~Li,
``Multi-probe LSH: Efficient indexing for high-dimensional similarity search,''
in \emph{Proc.\ 33rd International Conference on Very Large Data Bases (VLDB)},
pp.~950--961, 2007.

\bibitem{ailon2009fast}
N.~Ailon and B.~Chazelle,
``The fast Johnson--Lindenstrauss transform and approximate nearest neighbors,''
\emph{SIAM Journal on Computing}, vol.~39, no.~1, pp.~302--322, 2009.

\bibitem{jegou2011pq}
H.~J\'egou, M.~Douze, and C.~Schmid,
``Product quantization for nearest neighbor search,''
\emph{IEEE Transactions on Pattern Analysis and Machine Intelligence},
vol.~33, no.~1, pp.~117--128, 2011.

\bibitem{ge2013opq}
T.~Ge, K.~He, Q.~Ke, and J.~Sun,
``Optimized product quantization,''
\emph{IEEE Transactions on Pattern Analysis and Machine Intelligence},
vol.~36, no.~4, pp.~744--755, 2014.
(Conference version at CVPR 2013.)

\bibitem{malkov2014nsw}
Y.~A. Malkov, A.~Ponomarenko, A.~Logvinov, and V.~Krylov,
``Approximate nearest neighbor algorithm based on navigable small world graphs,''
\emph{Information Systems}, vol.~45, pp.~61--68, 2014.

\bibitem{bernhardsson2015annoy}
E.~Bernhardsson,
``Annoy: Approximate Nearest Neighbors in C++/Python,''
2015.
\url{https://github.com/spotify/annoy}.

\bibitem{malkov2018hnsw}
Y.~A. Malkov and D.~A. Yashunin,
``Efficient and robust approximate nearest neighbor search using
Hierarchical Navigable Small World graphs,''
\emph{IEEE Transactions on Pattern Analysis and Machine Intelligence},
vol.~42, no.~4, pp.~824--836, 2020.
(arXiv preprint 2016; journal version 2018/2020.)

\bibitem{johnson2019faiss}
J.~Johnson, M.~Douze, and H.~J\'egou,
``Billion-scale similarity search with GPUs,''
\emph{IEEE Transactions on Big Data}, vol.~7, no.~3, pp.~535--547, 2021.
(Faiss library; conference version at GTC 2017.)

\bibitem{jayaram2019diskann}
S.~Jayaram Subramanya, F.~Devvrit, H.~V. Simhadri, R.~Krishnaswamy,
and R.~Kadekodi,
``DiskANN: Fast accurate billion-point nearest neighbor search on a single node,''
in \emph{Advances in Neural Information Processing Systems (NeurIPS)},
vol.~32, 2019.

\bibitem{guo2020scann}
R.~Guo, P.~Sun, E.~Lindgren, Q.~Geng, D.~Simcha, F.~Chern, and S.~Kumar,
``Accelerating large-scale inference with anisotropic vector quantization,''
in \emph{Proc.\ 37th International Conference on Machine Learning (ICML)},
pp.~3887--3896, 2020.

\bibitem{khattab2020colbert}
O.~Khattab and M.~Zaharia,
``ColBERT: Efficient and effective passage search via contextualized late
interaction over BERT,''
in \emph{Proc.\ 43rd International ACM SIGIR Conference on Research and
Development in Information Retrieval}, pp.~39--48, 2020.

\bibitem{aumuller2020annbenchmarks}
M.~Aum\"uller, E.~Bernhardsson, and A.~Faithfull,
``ANN-Benchmarks: A benchmarking tool for approximate nearest neighbor algorithms,''
\emph{Information Systems}, vol.~87, 101374, 2020.

\bibitem{kusupati2022matryoshka}
A.~Kusupati, G.~Bhatt, A.~Rber, M.~Wallingford, A.~Sinha,
V.~Ramasamy, S.~N. Prasad, K.~Simhadri, and P.~Jain,
``Matryoshka representation learning,''
in \emph{Advances in Neural Information Processing Systems (NeurIPS)},
vol.~35, 2022.

\bibitem{babenko2014aq}
A.~Babenko and V.~Lempitsky,
``Additive quantization for extreme vector compression,''
in \emph{Proc.\ IEEE Conference on Computer Vision and Pattern Recognition
(CVPR)}, pp.~931--938, 2014.

\bibitem{chen2010rq}
Y.~Chen, T.~Guan, and C.~Wang,
``Approximate nearest neighbor search by residual vector quantization,''
\emph{Sensors}, vol.~10, no.~12, pp.~11259--11273, 2010.

\bibitem{dasgupta2008random}
S.~Dasgupta and Y.~Freund,
``Random projection trees and low dimensional manifolds,''
in \emph{Proc.\ 40th ACM Symposium on Theory of Computing (STOC)},
pp.~537--546, 2008.

\bibitem{bentley1975kdtree}
J.~L. Bentley,
``Multidimensional binary search trees used for associative searching,''
\emph{Communications of the ACM}, vol.~18, no.~9, pp.~509--517, 1975.

\bibitem{omohundro1989balltree}
S.~M. Omohundro,
``Five balltree construction algorithms,''
\emph{Technical Report TR-89-063}, International Computer Science
Institute, Berkeley, 1989.

\bibitem{douze2024faiss}
M.~Douze, A.~Guzhva, C.~Deng, J.~Johnson, G.~Szilvasy, P.-E.~Mazar\'e,
M.~Lomeli, L.~Hosseini, and H.~J\'egou,
``The Faiss library,''
\emph{arXiv preprint arXiv:2401.08281}, 2024.

\bibitem{aguerrebere2024similarity}
C.~Aguerrebere, I.~Bhatt, M.~Doherty, and R.~Jain,
``Similarity search in the blink of an eye with compressed indices,''
\emph{Proceedings of the VLDB Endowment}, vol.~16, no.~11, pp.~3433--3446,
2023.

\end{thebibliography}


\end{document}